\chapter{Discussion, Limitations, Synopsis Gap Closure, and Future Work}

\section{What the Thesis Has Actually Achieved}

The thesis has achieved a substantial transformation of the original CyclesGym base repository. It now supports Pakistan-adapted weather and soil inputs, Pakistan-oriented yearly price series, NPK-aware fertilization logic, hierarchical crop-planning integration, and structured reporting suitable for thesis writing. These are not superficial modifications. They change the scope, interpretability, and scientific defensibility of the system.

A core strength of the work is that it treats engineering as part of the research contribution. In simulator-based RL projects, progress is often described mainly through final returns. That is insufficient here. The thesis demonstrates that a credible agricultural RL platform also requires localized data provenance, explicit scope control, standardized outputs, and enough testing discipline that the final experiment campaign can be trusted when it runs.

Another achievement is conceptual clarification. The current repository now draws a sharper line between the role of the simulator, the role of RL, and the role of academic reporting. The simulator provides the transition dynamics. RL provides a way to optimize sequential decisions under those dynamics. The reporting layer provides the evidence needed to explain what the optimizer actually did. Treating these three roles separately has improved both the codebase and the thesis narrative.

\section{Why Scope Control is a Scientific Strength}

The most visible limitation of the thesis is also one of its strengths: it refuses to claim completion of everything promised in the synopsis. Irrigation as a learned action is not implemented. Rice-specific localization is not implemented. The main experimental crop setup remains maize-soy. The full broad latest-NPK matrix has been defined but not yet completed at build time. These facts narrow the current empirical chapter, but they also protect the thesis from making indefensible claims.

In many research settings, the temptation is to smooth over these gaps by writing a broader story than the code can support. This thesis takes the opposite approach. It uses a synopsis closure matrix to make the gaps visible and to distinguish clearly between completed engineering, partial implementation, and future work.

\input{tables/generated/synopsis_closure_matrix_nolabel.tex}

Scope control also has a practical defense benefit. It allows difficult questions to be answered directly. If a committee asks whether the thesis optimizes irrigation, the answer is no. If it asks whether the thesis completed a Pakistan-wide rice-wheat evaluation, the answer is no. If it asks whether the thesis nevertheless built a more localized, cost-aware, and more interpretable CyclesGym-derived platform, the answer is yes. This asymmetry is acceptable because a thesis is judged on defensible contribution, not on maximal ambition alone.

\section{Limitations}

The thesis has several technical and scientific limitations.
\begin{itemize}[leftmargin=1.2cm]
    \item The main broad latest-NPK campaign remains pending, so Chapter~6 is currently provisional.
    \item The active crop-planning and hierarchical experiments are still tied to the current working maize-soy configuration rather than a rice-wheat configuration.
    \item Irrigation has not yet been exposed as a learned control variable in the active training flows.
    \item The results remain simulation-based and have not yet been field-validated.
    \item Economic realism has improved substantially, but broader farm-level costs and market frictions are still simplified.
\end{itemize}

These limitations do not invalidate the completed work, but they define the boundary of what the thesis can responsibly claim.

\input{tables/generated/limitations_matrix.tex}

The limitations matrix is included because not all limitations have the same character. Some are evidence limitations, such as a pending experiment campaign. Others are modeling limitations, such as simplified economics. Others are scope limitations inherited from the fact that the thesis chose to stabilize one working crop configuration instead of spreading effort across many unstable configurations. Treating all limitations as one undifferentiated list would make later prioritization weaker.

\section{Implications for a Thesis Defense}

From a defense perspective, the thesis is strongest when it is presented as a system-construction and experimental-readiness contribution grounded in agricultural RL. It can credibly claim that the original base repository has been transformed into a localized and more interpretable research platform. It can also credibly explain why the fresh campaign must be executed before final comparative conclusions are frozen. This combination of technical depth and scope honesty is preferable to a superficially broader but less defensible narrative.

The historical benchmark context can still be useful in a defense because it shows that the repository is not a blank slate. Yet the defense should repeatedly emphasize that older runs are contextual and that the main thesis evidence is locked to the latest broad NPK design.

A disciplined defense posture therefore has two parts. First, it demonstrates mastery of the implemented system: environments, rewarders, pricing utilities, reporting outputs, and runner semantics. Second, it avoids claiming that simulator-based performance is equivalent to field-ready agronomy. This balance is likely to be viewed as maturity rather than weakness, especially in a technically demanding interdisciplinary topic.

\section{Future Work}

Future work follows naturally from the current limitations.
\begin{enumerate}[leftmargin=1.2cm]
    \item Complete the full broad latest-NPK matrix and regenerate all Chapter~6 figures from the completed summaries.
    \item Add irrigation as a learned action so that the thesis scope aligns more closely with the original synopsis.
    \item Retarget the crop-planning stack from the current working maize-soy setup toward a Pakistan-focused cereal rotation if agronomic data and configuration stability permit.
    \item Add stronger statistical analysis, including confidence intervals and formal significance testing across seeds where coverage is sufficient.
    \item Conduct expert-in-the-loop or field-adjacent validation to test whether the learned policies remain plausible outside the simulator.
\end{enumerate}

\begin{figure}[H]
    \centering
    \includegraphics[width=0.90\textwidth]{future_work_roadmap.pdf}
    \caption{Roadmap implied by the current thesis limitations. The roadmap is staged so that empirical completion, control expansion, and external validation occur in a defensible order.}
    \label{fig:future-work-roadmap}
\end{figure}

The order shown in Figure~\ref{fig:future-work-roadmap} is intentional. The thesis should not jump to new crops, irrigation, or validation studies before completing the already committed broad NPK campaign on the stabilized working stack. Otherwise, the project risks accumulating new branches of scope faster than it closes the main one.

\section{Final Reflection}

The central message of this thesis is not that agricultural RL is solved. It is that a vague idea of cost-aware crop management has been converted into a concrete, localized, testable, and reportable software stack. Even before the final broad campaign completes, that transformation has scientific value. It turns an inherited simulator wrapper into a research platform with clearer provenance, better scope control, and a more rigorous path to defensible experimentation. Once the fresh broad campaign is executed, the thesis can add the full comparative chapter on top of this foundation without needing to rewrite the underlying story.
